El efecto fotoeléctrico constituye uno de los experimentos fundamentales de la
física moderna, pues demuestra que la energía de la luz está cuantizada en
fotones de energía $E=h\nu$, donde $h$ es la constante de Planck y $\nu$ la
frecuencia de la radiación. La explicación propuesta por Albert Einstein en
1905 permitió interpretar diversos fenómenos que no podían describirse mediante
la teoría clásica del electromagnetismo y representó uno de los primeros
evidencias de la naturaleza cuántica de la luz. \\

Un diodo emisor de luz (LED) funciona mediante la recombinación de electrones y
huecos en un material semiconductor. Durante este proceso se emiten fotones cuya
energía está relacionada con la diferencia de energía entre las bandas del
material. Como primera aproximación puede considerarse que la energía adquirida
por un electrón al atravesar la unión del diodo es igual a la energía del fotón
emitido,

\[
eV=h\nu,
\]

donde $e$ es la carga elemental y $V$ es el voltaje umbral del LED. Esta
expresión predice una relación lineal entre la frecuencia de la luz emitida y el
voltaje aplicado, cuya pendiente corresponde a $e/h$.\\

En esta práctica se utilizaron LEDs de diferentes longitudes de onda para medir
su voltaje de umbral y construir la gráfica de frecuencia contra voltaje.
Mediante un ajuste lineal se obtuvo una estimación experimental de la relación
$h/e$, comparándola posteriormente con el valor aceptado y evaluando el error
porcentual del método.